\documentclass{article}
\usepackage{mathtools} 
\usepackage{fontspec}
\usepackage[UTF8]{ctex}
\usepackage{amsthm}
\usepackage{mdframed}
\usepackage{xcolor}
\usepackage{amssymb}
\usepackage{amsmath}


% 定义新的带灰色背景的说明环境 zremark
\newmdtheoremenv[
  backgroundcolor=gray!10,
  % 边框与背景一致，边框线会消失
  linecolor=gray!10
]{zremark}{说明}

% 通用矩阵命令: \flexmatrix{矩阵名}{元素符号}{行数}{列数}
\newcommand{\flexmatrix}[4]{
  \[
  #1 = \begin{pmatrix}
    #2_{11}     & #2_{12}     & \cdots & #2_{1#4}   \\
    #2_{21}     & #2_{22}     & \cdots & #2_{2#4}   \\
    \vdots      & \vdots      & \ddots & \vdots     \\
    #2_{#31}    & #2_{#32}    & \cdots & #2_{#3#4}
  \end{pmatrix}
  \]
}

% 简化版命令（默认矩阵名为A，元素符号为a）: \quickmatrix{行数}{列数}
\newcommand{\quickmatrix}[2]{\flexmatrix{A}{a}{#1}{#2}}

\begin{document}
\title{3.3 注释}
\author{张志聪}
\maketitle

\begin{zremark}
  归结原则（$x \to +\infty$类型极限）：
  设$f$在$U(+\infty)$上有定义。$\lim\limits_{x \to +\infty} f(x)$存在的充分必要条件是：
  任何以$+\infty$为极限的数列$\{x_n\}$，极限$\lim\limits_{n \to +\infty} f(x_n)$存在且相等。
\end{zremark}

\textbf{证明：}

\begin{itemize}
  \item 必要性

        设$\lim\limits_{x \to +\infty} f(x) = A$，
        则对任给的$\epsilon > 0$，存在$M > 0$，使得当$x > M$时，有
        \begin{align*}
          |f(x) - A| < \epsilon
        \end{align*}
        另一方面，设数列$\{x_n\}$且$\lim\limits_{n \to +\infty} x_n = +\infty$，
        则对上述的$M$，存在$N$，使得$n > N$时，有
        \begin{align*}
          x_n > M
        \end{align*}
        从而有
        \begin{align*}
          |f(x_n) - A| < \epsilon
        \end{align*}
        所以$\lim\limits_{n \to +\infty} f(x_n) = A$。

  \item 充分性

        设对任何数列$\{x_n\}$且$\lim\limits_{n \to +\infty} x_n = +\infty$，有
        $\lim\limits_{n \to +\infty} f(x_n) = A$，则可用反证法推出
        $\lim\limits_{x \to +\infty} f(x) = A$。
        事实上，倘若当$x \to \infty$时$f$不以$A$为极限，则存在某个$\epsilon_0 > 0$，
        对任何$M > 0$（不论多么大），总存在一点$x$，尽管$x > M$，但有
        $|f(x) - A| \geq \epsilon_0$。现依次取$M = 1, 2, \cdots, n, \cdots$，
        则存在相应的点$x_1, x_2, \cdots, x_n, \cdots$，使得
        \begin{align*}
          \lim\limits_{x \to \infty} x_n = +\infty
        \end{align*}
        但当$n \to +\infty$时$f(x_n)$不趋于$A$，故存在矛盾，假设不成立，命题得证。
\end{itemize}


\end{document}